
\section{Physical conditions of renormalisation}

\SourcePage{538}{543}

The theory presented so far in this chapter has been to a significant extent formal in character. We operated with all quantities as if they were finite, intentionally leaving aside the infinities arising in the theory. However, in an actual computation of the functions~$\mathcal{D}$, $\mathcal{G}$, $\Gamma$ by perturbation theory, divergent integrals appear which cannot, without additional considerations, be ascribed any definite value. The emergence of such divergences is a manifestation of the logical incompleteness of quantum electrodynamics. We shall see, however, that in this theory one can establish certain prescriptions that allow an unambiguous ``subtraction of infinities'' and the obtaining of finite values for all quantities possessing direct physical meaning. At the basis of these prescriptions lie obvious physical requirements: that the photon mass be zero and that the charge and mass of the electron equal their observed values.

Let us begin with the conditions on the photon propagator.

Consider scattering through single-particle intermediate states containing one virtual photon. The amplitude must have a pole at~$P^2 = 0$ (the square of the total 4-momentum of the initial particles~$P$ coincides with the square of the real photon mass). As we saw in~\S79, this follows from the unitarity condition. The pole term arises from diagrams of the type~(79.1):

% [Feynman diagram (110.1): scattering via single virtual photon --- two shaded blobs (representing exact scattering amplitudes) connected by a bold dashed line (exact photon propagator) carrying momentum k = P]
\begin{equation}
\label{eq:110.1}
\end{equation}

\noindent with radiative corrections, both parts of the diagram being connected by a bold dashed line (exact photon propagator). Therefore~$\mathcal{D}(k^2)$ must have a pole at~$k^2 = 0$:
\begin{equation}
\mathcal{D} \to \frac{4\pi Z}{k^2} \quad \text{when} \quad k^2 \to 0,
\label{eq:110.2}
\end{equation}
where~$Z$ is a constant. For the polarisation operator this means, according to~(103.21):
\begin{equation}
\mathcal{P}(0) = 0.
\label{eq:110.3}
\end{equation}
The coefficient~$1/Z = 1 - \mathcal{P}(k^2)/k^2\big|_{k^2 \to 0}$.

Further restrictions on~$\mathcal{P}(k^2)$ follow from the physical definition of electric charge. Two classical (arbitrarily heavy) particles at rest at a large distance from one another must interact according to the Coulomb law~$U = e^2/r$ ($r \gg 1/m$). On the other hand, this interaction is expressed by the diagram:

% [Feynman diagram (110.4): scattering of two heavy classical particles via a single virtual photon --- upper and lower straight lines (heavy particles a,b and c,d) connected by a bold dashed photon line carrying momentum k]
\begin{equation}
\label{eq:110.4}
\end{equation}

\noindent where the upper and lower lines correspond to the classical particles. The photon self-energy corrections are included in the virtual photon line. Other corrections, involving the heavy-particle lines, vanish in the limit~$M \to \infty$: adding internal lines to the diagram~(\ref{eq:110.4}) (for example, connecting lines~$a$ and~$c$, or~$a$ and~$b$, by a photon line) leads to virtual heavy-particle propagators containing the mass~$M$ in the denominator and vanishing as~$M \to \infty$.

\SourcePage{539}{544}

From the diagram~(\ref{eq:110.4}), the factor~$e^2\mathcal{D}(k^2)$ must represent (up to sign) the Fourier transform of the interaction potential. Stationarity of the particles means the virtual photon frequency~$\omega = 0$; large distances mean small wave vectors~$\mathbf{k}$. The Fourier transform of the Coulomb potential is~$4\pi e^2/\mathbf{k}^2$. Since~$\mathcal{D}$ depends only on~$k^2 = \omega^2 - \mathbf{k}^2$, we obtain:
\begin{equation}
\mathcal{D} \to \frac{4\pi}{k^2}, \quad k^2 \to 0.
\label{eq:110.5}
\end{equation}
Thus the coefficient~$Z$ in~(\ref{eq:110.2}) must equal unity: $Z = 1$ (the sign in~(\ref{eq:110.5}) is obvious: $\mathcal{D}(k^2)$ tends to the free photon propagator~$D(k^2)$). For the polarisation operator this means:
\begin{equation}
\frac{\mathcal{P}(k^2)}{k^2} \to 0, \quad k^2 \to 0.
\label{eq:110.6}
\end{equation}
Besides condition~(\ref{eq:110.3}), this also gives:
\begin{equation}
\mathcal{P}'(0) = 0.
\label{eq:110.7}
\end{equation}

In~\S103 it was noted that the effective external photon line factor~(103.15), with account of~(103.16) and~(103.20), is:
\[
\Bigl[1 + \frac{\mathcal{P}(0)\,\mathcal{D}(0)}{4\pi}\Bigr]\,e^\mu.
\]
Now, with~(\ref{eq:110.5}) and~(\ref{eq:110.6}), the correction term vanishes. In other words, in external photon lines one need not include radiative corrections at all.

Thus natural physical requirements lead to the establishment of definite (zero) values of~$\mathcal{P}(0)$ and~$\mathcal{P}'(0)$. Computing these quantities from perturbation-theory diagrams would lead to divergent integrals. The method of eliminating infinities consists in ascribing to the divergent expressions the values prescribed by the physical requirements. This procedure is called \emph{renormalisation} of the corresponding quantities.\footnote{The idea of this approach was first expressed by Kramers (H.~Kramers, 1947). Systematic application of the renormalisation method in quantum electrodynamics was carried out in the works of Dyson, Tomonaga (S.~Tomonaga), Feynman and Schwinger.}

\SourcePage{540}{545}

The procedure can also be formulated differently. For charge renormalisation one introduces a non-physical ``bare'' charge~$e_0$ as the parameter in the electromagnetic interaction Hamiltonian in the formal perturbation theory. After this, the renormalisation condition is formulated as: $e_0^2\,\mathcal{D}(k^2) \to 4\pi e^2/k^2$ (at~$k^2 \to 0$), where~$e$ is the true physical charge. From this~$e_0^2\,Z = e^2$, and together with~$Z$ the non-physical quantity~$e_0$ is eliminated. By requiring~$Z = 1$ from the outset, we perform the renormalisation ``on the fly'' and avoid introducing fictitious quantities even in intermediate calculations.

Now let us turn to the conditions on the electron propagator. Consider scattering through a single-particle state containing one virtual electron. The amplitude must have a pole at~$P^2 = m^2$ (the square of the total 4-momentum of the initial particles~$P$ equals the square of the real electron mass).
